% Generated draft: requires mathematical and source review.
\documentclass[11pt]{article}
\usepackage[margin=25mm]{geometry}
\usepackage{fontspec}
\setmainfont{Latin Modern Roman}
\usepackage{amsmath,amssymb,mathtools,unicode-math}
\setmathfont{Latin Modern Math}
\usepackage{enumitem,needspace,xurl,hyperref}
\setlength{\parindent}{0pt}
\setlength{\parskip}{5pt}
\setlength{\emergencystretch}{3em}
\newcommand{\N}{\mathbb{N}}
\newcommand{\Z}{\mathbb{Z}}
\newcommand{\Q}{\mathbb{Q}}
\newcommand{\R}{\mathbb{R}}
\newcommand{\C}{\mathbb{C}}
\newcommand{\F}{\mathbb{F}}
\newcommand{\A}{\mathbb{A}}
\newcommand{\PP}{\mathbb{P}}
\newcommand{\E}{\mathbb{E}}
\newcommand{\eps}{\varepsilon}
\newcommand{\dd}{\,\mathrm{d}}
\begin{document}
\section*{UnsolvedMath ID: 4000006}
\textbf{Source problem code:} catalogue entry
\subsection{Definitions and mathematical statement}
The target is a method for recovering a known sharp theorem. In a standard formulation, let $(M,g)$ be a connected closed Riemannian manifold of dimension $N\geq2$, and suppose
\[
\operatorname{Ric}_x(v,v)\geq\rho\,g_x(v,v)\quad(x\in M,\ v\in T_xM),\qquad \rho>0.
\]
Let $\lambda_1(-\Delta_g)$ be the first positive eigenvalue of the nonnegative Laplace--Beltrami operator. The sharp Lichnerowicz estimate is
\[
\lambda_1(-\Delta_g)\geq\frac{N}{N-1}\rho,
\]
whereas the source says the coarse-curvature argument recovers only the scale $\lambda_1\geq\rho$. Here $\inf\operatorname{Ric}$ in the source is understood as the infimum over unit tangent vectors, not an infimum of an uncontracted tensor.

For the source's $\varepsilon$-step random walk, denote its transition law from $x$ by $P_x^\varepsilon$. Coarse curvature uses
\[
\kappa_\varepsilon(x,y)=1-\frac{W_1(P_x^\varepsilon,P_y^\varepsilon)}{d_g(x,y)},\qquad x\ne y,
\]
where $W_1$ minimizes expected geodesic distance over couplings of the two laws. The problem asks whether one can retain the directional information lost by this scalar contraction argument, for example using reflection couplings, to obtain the full factor $N/(N-1)$ in the Lichnerowicz conclusion.

Closedness, connectedness, positivity of $\rho$, and $N\geq2$ make the conventional displayed spectral theorem precise; these hypotheses are not all supplied in the source. The step distribution (such as a geodesic sphere or ball), scaling of the discrete spectral gap to the Laplacian, and admissible curvature/coupling refinement are also not specified. The question concerns deriving the sharp estimate by the indicated method, not whether the classical estimate is true.
\subsection{Short English statement}
Can a refinement of coarse Ricci curvature, perhaps through reflection couplings that retain direction, recover the sharp $N/(N-1)$ factor in the Lichnerowicz spectral-gap theorem? The theorem is the known target; the requested advance is the coupling or curvature argument.
\subsection{Sources}
\begin{itemize}
\item Local UnsolvedMath registry, permanent ID 4000006, exact source code catalogue entry, ``Sharp Lichnerowicz theorem''; proposed by Yann Ollivier, 2008. Registry status partially solved, classification PARTIAL-PROGRESS.
\item Yann Ollivier, \emph{Discrete Ricci curvature: Open problems} (2008), Problem F, \url{https://www.yann-ollivier.org/rech/publs/problems_curvmarkov.pdf}. Supplied pdf-text extraction accessed July 29, 2026; the document presents the item as open, with no later authoritative resolution check recorded. No PDF was accessed for this draft.
\item The supplied summary reports improvements using reflection couplings and directional structure, while saying recovery of the sharp constant by a coarse or multistep curvature argument remains open. It gives no precise later theorem or reference, so this is retained as a supplied partial-progress assessment.
\item Provenance flags status and rights NEEDS\_REVIEW: public access was verified, but redistribution permission was not established.
\end{itemize}
\end{document}
